% Fichier engendré par tools/generate-tex.py à partir de 01-nombres-et-calculs.
% Ne pas éditer à la main : tout le texte provient des commentaires du fichier Lean.
\documentclass[11pt,a4paper]{article}

% Compile aussi bien avec pdflatex qu'avec xelatex ou tectonic.
\usepackage{iftex}
\ifPDFTeX
  \usepackage[T1]{fontenc}
  \usepackage[utf8]{inputenc}
\else
  \usepackage{fontspec}
\fi
\usepackage[french]{babel}
\usepackage{amsmath,amssymb,amsthm}
\usepackage[margin=2.5cm]{geometry}
\usepackage[hidelinks]{hyperref}

\theoremstyle{definition}
\newtheorem{definition}{Définition}
\newtheorem{exemple}{Exemple}
\theoremstyle{plain}
\newtheorem{theoreme}{Théorème}
\newtheorem{lemme}[theoreme]{Lemme}

% Un exemple ne se démontre pas : on explique ce qu'il montre.
\newenvironment{explication}{\begin{proof}[Explication]}{\end{proof}}

% Renvoi à la déclaration Lean, une ligne après la démonstration, aligné à droite.
\newcommand{\source}[2]{\par\smallskip\noindent\hfill{\footnotesize\href{#1}{\texttt{#2}}}\par\smallskip}

\title{Nombres et calculs}
\date{}

\begin{document}
\maketitle
% Les identifiants Lean en machine à écrire ne se coupent pas : on tolère des
% espacements plus lâches plutôt que des lignes qui débordent.
\sloppy

\section{CalculsEtCalculLitteral}

\noindent
Collège \textemdash{} section \og{} Calculs et calcul littéral \fg{}.
Énoncés et démonstrations en français : voir NombresEtCalculs.tex.

\subsection{Priorités opératoires ; rôle des parenthèses}

\begin{exemple}
La multiplication est prioritaire sur l'addition : sans parenthèses, \texttt{2 + 3 \(\times\) 4}
vaut \texttt{14}, et non \texttt{20}. La priorité est une convention d'écriture, pas un théorème ;
ce qui se démontre, c'est que les deux expressions diffèrent.
\end{exemple}

\begin{explication}
Les deux calculs sont menés à leur terme et comparés : $2 + 3 \times 4 = 14$ tandis
que $(2+3) \times 4 = 20$. La priorité de la multiplication est une convention d'écriture,
pas un théorème ; ce qui se démontre, c'est que les deux résultats diffèrent.
\end{explication}
\source{https://github.com/Commutator-IO/learn-lean/blob/main/courses/01-college/01-nombres-et-calculs/CalculsEtCalculLitteral.lean\#L14}{CalculsEtCalculLitteral.lean\#L14}

\subsection{Addition et soustraction des relatifs ; \texttt{a − b = a + (−b)}}

\begin{theoreme}
Soustraire, c'est ajouter l'opposé.
\end{theoreme}

\begin{proof}
C'est la définition de la soustraction dans les relatifs : $a - b$ est posé comme
$a + (-b)$.
\end{proof}
\source{https://github.com/Commutator-IO/learn-lean/blob/main/courses/01-college/01-nombres-et-calculs/CalculsEtCalculLitteral.lean\#L20}{CalculsEtCalculLitteral.lean\#L20}

\begin{theoreme}
L'opposé d'une somme est la somme des opposés.
\end{theoreme}

\begin{proof}
L'opposé d'une somme s'obtient en prenant l'opposé de chaque terme, propriété du
groupe additif des entiers relatifs.
\end{proof}
\source{https://github.com/Commutator-IO/learn-lean/blob/main/courses/01-college/01-nombres-et-calculs/CalculsEtCalculLitteral.lean\#L23}{CalculsEtCalculLitteral.lean\#L23}

\subsection{Règle des signes pour la multiplication et la division}

\begin{theoreme}
Le produit de deux nombres de signes contraires est négatif.
\end{theoreme}

\begin{proof}
Changer le signe d'un facteur change le signe du produit : $(-a) \times b$ et
$a \times b$ sont opposés.
\end{proof}
\source{https://github.com/Commutator-IO/learn-lean/blob/main/courses/01-college/01-nombres-et-calculs/CalculsEtCalculLitteral.lean\#L28}{CalculsEtCalculLitteral.lean\#L28}

\begin{theoreme}
Le produit de deux nombres négatifs est positif.
\end{theoreme}

\begin{proof}
En appliquant deux fois la règle précédente : $(-a)(-b) = -(a(-b)) = -(-(ab)) = ab$.
\end{proof}
\source{https://github.com/Commutator-IO/learn-lean/blob/main/courses/01-college/01-nombres-et-calculs/CalculsEtCalculLitteral.lean\#L31}{CalculsEtCalculLitteral.lean\#L31}

\begin{theoreme}
Même règle pour le quotient.
\end{theoreme}

\begin{proof}
Un quotient est un produit par l'inverse, la règle des signes s'y transporte donc
telle quelle.
\end{proof}
\source{https://github.com/Commutator-IO/learn-lean/blob/main/courses/01-college/01-nombres-et-calculs/CalculsEtCalculLitteral.lean\#L34}{CalculsEtCalculLitteral.lean\#L34}

\subsection{Somme, produit, quotient de fractions ; diviser = multiplier par l'inverse}

\begin{theoreme}
Somme de deux fractions.
\end{theoreme}

\begin{proof}
On réduit au même dénominateur $bd$, non nul puisque $b$ et $d$ le sont :
\[
\frac{a}{b} + \frac{c}{d} = \frac{ad}{bd} + \frac{cb}{bd} = \frac{ad + cb}{bd} .
\]
\end{proof}
\source{https://github.com/Commutator-IO/learn-lean/blob/main/courses/01-college/01-nombres-et-calculs/CalculsEtCalculLitteral.lean\#L39}{CalculsEtCalculLitteral.lean\#L39}

\begin{theoreme}
Produit de deux fractions.
\end{theoreme}

\begin{proof}
Les quotients se multiplient numérateurs entre eux et dénominateurs entre eux :
\[
\frac{a}{b} \times \frac{c}{d}
  = a \times \frac1b \times c \times \frac1d
  = (ac) \times \frac{1}{bd}
  = \frac{ac}{bd} .
\]
\end{proof}
\source{https://github.com/Commutator-IO/learn-lean/blob/main/courses/01-college/01-nombres-et-calculs/CalculsEtCalculLitteral.lean\#L44}{CalculsEtCalculLitteral.lean\#L44}

\begin{theoreme}
Diviser par une fraction, c'est multiplier par son inverse.
\end{theoreme}

\begin{proof}
On écrit les deux divisions comme des produits par l'inverse, puis on réordonne les
facteurs :
\[
\frac{a/b}{c/d} = \frac{a}{b} \times \frac{1}{c/d}
                = \frac{a}{b} \times \frac{d}{c} .
\]

L'énoncé mathématique suppose $c \ne 0$ et $d \ne 0$ : sans cela, aucun des deux membres
n'a de sens, la division par zéro n'étant pas définie.

Le théorème Lean, lui, est énoncé sans ces hypothèses, et c'est un artefact de la
formalisation : Lean exige que toute fonction soit totale, et Mathlib pose donc
$x/0 = 0$ pour compléter la division. Les deux membres valent alors zéro, et l'égalité
est vraie pour de mauvaises raisons. Cette valeur de remplissage ne dit rien des
mathématiques : elle rend seulement la fonction définie partout.
\end{proof}
\source{https://github.com/Commutator-IO/learn-lean/blob/main/courses/01-college/01-nombres-et-calculs/CalculsEtCalculLitteral.lean\#L48}{CalculsEtCalculLitteral.lean\#L48}

\subsection{Distributivité simple : \texttt{k(a + b) = ka + kb}}

\begin{theoreme}
Distributivité de la multiplication sur l'addition.
\end{theoreme}

\begin{proof}
C'est la distributivité de la multiplication sur l'addition.
\end{proof}
\source{https://github.com/Commutator-IO/learn-lean/blob/main/courses/01-college/01-nombres-et-calculs/CalculsEtCalculLitteral.lean\#L55}{CalculsEtCalculLitteral.lean\#L55}

\subsection{Double distributivité : \texttt{(a + b)(c + d) = ac + ad + bc + bd}}

\begin{theoreme}
Développement d'un produit de deux sommes.
\end{theoreme}

\begin{proof}
On développe deux fois :
\[
(a+b)(c+d) = a(c+d) + b(c+d) = ac + ad + bc + bd .
\]
\end{proof}
\source{https://github.com/Commutator-IO/learn-lean/blob/main/courses/01-college/01-nombres-et-calculs/CalculsEtCalculLitteral.lean\#L60}{CalculsEtCalculLitteral.lean\#L60}

\subsection{Factorisation d'une expression à facteur commun}

\begin{theoreme}
Factoriser, c'est lire la distributivité de droite à gauche.
\end{theoreme}

\begin{proof}
C'est la distributivité lue de droite à gauche : le membre de gauche est le
développement du membre de droite.
\end{proof}
\source{https://github.com/Commutator-IO/learn-lean/blob/main/courses/01-college/01-nombres-et-calculs/CalculsEtCalculLitteral.lean\#L66}{CalculsEtCalculLitteral.lean\#L66}

\subsection{Un produit est nul \(\iff\) l'un des facteurs est nul}

\begin{theoreme}
L'équation produit : un produit de deux nombres est nul si et seulement si l'un des
deux est nul.
\end{theoreme}

\begin{proof}
Si l'un des facteurs est nul, le produit l'est. Réciproquement, si $a \ne 0$, on peut
diviser $ab = 0$ par $a$, ce qui donne $b = 0$.
\end{proof}
\source{https://github.com/Commutator-IO/learn-lean/blob/main/courses/01-college/01-nombres-et-calculs/CalculsEtCalculLitteral.lean\#L73}{CalculsEtCalculLitteral.lean\#L73}

\subsection{Puissances : \texttt{aᵐ \(\times\) aⁿ = aᵐ⁺ⁿ}, \texttt{aᵐ / aⁿ = aᵐ⁻ⁿ}, \texttt{(aᵐ)ⁿ = aᵐⁿ}, \texttt{(ab)ⁿ = aⁿbⁿ}}

\begin{theoreme}
Produit de deux puissances d'un même nombre.
\end{theoreme}

\begin{proof}
En comptant les facteurs :
\[
a^m \times a^n
  = \underbrace{a \times \cdots \times a}_{m}
    \times \underbrace{a \times \cdots \times a}_{n}
  = \underbrace{a \times \cdots \times a}_{m+n}
  = a^{m+n} .
\]
\end{proof}
\source{https://github.com/Commutator-IO/learn-lean/blob/main/courses/01-college/01-nombres-et-calculs/CalculsEtCalculLitteral.lean\#L78}{CalculsEtCalculLitteral.lean\#L78}

\begin{theoreme}
Quotient de deux puissances d'un même nombre non nul.
\end{theoreme}

\begin{proof}
Même décompte, avec des exposants entiers relatifs :
\[
\frac{a^m}{a^n} = a^m \times a^{-n} = a^{m-n} .
\]
L'hypothèse $a \ne 0$ est indispensable pour donner un sens aux exposants négatifs.
\end{proof}
\source{https://github.com/Commutator-IO/learn-lean/blob/main/courses/01-college/01-nombres-et-calculs/CalculsEtCalculLitteral.lean\#L82}{CalculsEtCalculLitteral.lean\#L82}

\begin{theoreme}
Puissance d'une puissance.
\end{theoreme}

\begin{proof}
$(a^m)^n$ est le produit de $n$ facteurs valant chacun $a^m$ :
\[
(a^m)^n = \underbrace{a^m \times \cdots \times a^m}_{n}
        = a^{\,\underbrace{\scriptstyle m + \cdots + m}_{n}} = a^{mn} .
\]
\end{proof}
\source{https://github.com/Commutator-IO/learn-lean/blob/main/courses/01-college/01-nombres-et-calculs/CalculsEtCalculLitteral.lean\#L86}{CalculsEtCalculLitteral.lean\#L86}

\begin{theoreme}
Puissance d'un produit.
\end{theoreme}

\begin{proof}
On regroupe, ce que permet la commutativité, les $n$ facteurs $a$ d'un côté et les $n$
facteurs $b$ de l'autre :
\[
(ab)^n = \underbrace{(ab) \times \cdots \times (ab)}_{n}
       = \underbrace{a \times \cdots \times a}_{n} \times
         \underbrace{b \times \cdots \times b}_{n}
       = a^n b^n .
\]
\end{proof}
\source{https://github.com/Commutator-IO/learn-lean/blob/main/courses/01-college/01-nombres-et-calculs/CalculsEtCalculLitteral.lean\#L90}{CalculsEtCalculLitteral.lean\#L90}

\subsection{\texttt{a⁻ⁿ = 1/aⁿ} pour \texttt{a \(\ne\) 0} ; \texttt{a⁰ = 1}}

\begin{theoreme}
Exposant négatif.
\end{theoreme}

\begin{proof}
Par définition de l'exposant négatif : $a^{-n}$ est l'inverse de $a^n$.
\end{proof}
\source{https://github.com/Commutator-IO/learn-lean/blob/main/courses/01-college/01-nombres-et-calculs/CalculsEtCalculLitteral.lean\#L96}{CalculsEtCalculLitteral.lean\#L96}

\begin{theoreme}
Exposant nul : $a^0 = 1$, que l'exposant soit un entier naturel ou un entier relatif.
\end{theoreme}

\begin{proof}
C'est une convention, et non un théorème : on pose $a^0 = 1$ pour que la règle
$a^m \times a^n = a^{m+n}$ vaille encore lorsque l'un des exposants est nul.

Si l'énoncé formel porte deux égalités et non une, c'est que Lean distingue deux
opérations là où l'écriture n'en montre qu'une. La puissance d'exposant entier naturel est
une multiplication répétée ; celle d'exposant entier relatif passe par l'inverse et se
définit à part. La convention doit donc être posée pour chacune. La distinction est celle
de la bibliothèque, pas celle des mathématiques.
\end{proof}
\source{https://github.com/Commutator-IO/learn-lean/blob/main/courses/01-college/01-nombres-et-calculs/CalculsEtCalculLitteral.lean\#L103}{CalculsEtCalculLitteral.lean\#L103}

\subsection{Conservation des inégalités}

\begin{theoreme}
Ajouter un même nombre aux deux membres conserve l'inégalité.
\end{theoreme}

\begin{proof}
Ajouter $c$ aux deux membres ne change pas leur différence :
\[
(b + c) - (a + c) = b - a \ge 0 .
\]
\end{proof}
\source{https://github.com/Commutator-IO/learn-lean/blob/main/courses/01-college/01-nombres-et-calculs/CalculsEtCalculLitteral.lean\#L109}{CalculsEtCalculLitteral.lean\#L109}

\begin{theoreme}
Multiplier par un nombre strictement positif conserve l'inégalité.
\end{theoreme}

\begin{proof}
Multiplier par $c$ multiplie la différence par $c$ :
\[
bc - ac = (b - a)c \ge 0 \quad\text{car } b - a \ge 0 \text{ et } c > 0 .
\]
\end{proof}
\source{https://github.com/Commutator-IO/learn-lean/blob/main/courses/01-college/01-nombres-et-calculs/CalculsEtCalculLitteral.lean\#L113}{CalculsEtCalculLitteral.lean\#L113}

\begin{theoreme}
Multiplier par un nombre strictement négatif renverse l'inégalité.
\end{theoreme}

\begin{proof}
Le même produit change de signe :
\[
ac - bc = (b - a)(-c) \ge 0 \quad\text{car } b - a \ge 0 \text{ et } -c > 0 ,
\]
d'où $bc \le ac$ : l'inégalité se renverse.
\end{proof}
\source{https://github.com/Commutator-IO/learn-lean/blob/main/courses/01-college/01-nombres-et-calculs/CalculsEtCalculLitteral.lean\#L117}{CalculsEtCalculLitteral.lean\#L117}

\subsection{Équation du premier degré \texttt{ax + b = 0} : solution unique si \texttt{a \(\ne\) 0}}

\begin{theoreme}
Une équation du premier degré à coefficient directeur non nul a une solution, et une
seule : \texttt{x = -b/a}.
\end{theoreme}

\begin{proof}
\emph{Existence.} $x = -b/a$ convient : $a \times (-b/a) + b = -b + b = 0$, le calcul
étant licite puisque $a \ne 0$.

\emph{Unicité.} Si $y$ est solution, alors $ay = -b = a \times (-b/a)$, et comme $a \ne 0$
on peut simplifier par $a$, d'où $y = -b/a$.
\end{proof}
\source{https://github.com/Commutator-IO/learn-lean/blob/main/courses/01-college/01-nombres-et-calculs/CalculsEtCalculLitteral.lean\#L124}{CalculsEtCalculLitteral.lean\#L124}

\subsection{Équation produit \texttt{(ax + b)(cx + d) = 0}}

\begin{theoreme}
Une équation produit se ramène à deux équations du premier degré.
\end{theoreme}

\begin{proof}
C'est l'équation produit appliquée aux deux facteurs $ax + b$ et $cx + d$ : le
produit est nul exactement quand l'un des deux l'est. Résoudre l'équation revient donc à
résoudre deux équations du premier degré.
\end{proof}
\source{https://github.com/Commutator-IO/learn-lean/blob/main/courses/01-college/01-nombres-et-calculs/CalculsEtCalculLitteral.lean\#L132}{CalculsEtCalculLitteral.lean\#L132}

\subsection{Tester si un nombre est solution ; démontrer qu'une égalité est vraie pour tout \texttt{x}}

\begin{exemple}
Tester une valeur : \texttt{2} est solution de \texttt{x² − 5x + 6 = 0}, \texttt{1} ne l'est pas.
\end{exemple}

\begin{explication}
Les deux calculs sont menés : $2^2 - 5 \times 2 + 6 = 0$, donc $2$ est solution ;
$1^2 - 5 \times 1 + 6 = 2 \ne 0$, donc $1$ ne l'est pas.
\end{explication}
\source{https://github.com/Commutator-IO/learn-lean/blob/main/courses/01-college/01-nombres-et-calculs/CalculsEtCalculLitteral.lean\#L138}{CalculsEtCalculLitteral.lean\#L138}

\begin{theoreme}
Démontrer une identité, c'est l'établir pour tout \texttt{x} à la fois \textemdash{} ce qu'aucun test de
valeurs ne remplace.
\end{theoreme}

\begin{proof}
Le développement du membre de gauche donne le membre de droite, quelle que soit la
valeur de $x$. C'est ce « quelle que soit » qui fait la démonstration, et qu'aucun test de
valeurs particulières ne remplace.
\end{proof}
\source{https://github.com/Commutator-IO/learn-lean/blob/main/courses/01-college/01-nombres-et-calculs/CalculsEtCalculLitteral.lean\#L143}{CalculsEtCalculLitteral.lean\#L143}

\subsection{Programme de calcul : deux programmes donnent le même résultat pour toute entrée}

\begin{theoreme}
\og{} Ajouter 1, élever au carré, retrancher 1 \fg{} et \og{} multiplier par le nombre augmenté
de 2 \fg{} donnent toujours le même résultat.
\end{theoreme}

\begin{proof}
En développant : $(x+1)^2 - 1 = x^2 + 2x + 1 - 1 = x^2 + 2x = x(x+2)$. Les deux
programmes coïncident donc pour toute entrée.
\end{proof}
\source{https://github.com/Commutator-IO/learn-lean/blob/main/courses/01-college/01-nombres-et-calculs/CalculsEtCalculLitteral.lean\#L149}{CalculsEtCalculLitteral.lean\#L149}

\section{EcrituresDesNombres}

\noindent
Collège \textemdash{} section \og{} Écritures des nombres \fg{}. Mathlib est nécessaire ici : les réels,
la racine carrée et l'irrationalité de \(\sqrt{2}\) n'existent pas dans Lean core.
Énoncés et démonstrations en français : voir NombresEtCalculs.tex.

\subsection{Numération décimale de position ; valeur d'un chiffre selon son rang}

\begin{theoreme}
Un entier vaut la somme de ses chiffres pondérés par les puissances de dix.
\end{theoreme}

\begin{proof}
C'est la propriété fondamentale de l'écriture décimale : reconstruire un nombre à
partir de la liste de ses chiffres, en les pondérant par les puissances successives de dix,
redonne le nombre de départ.
\end{proof}
\source{https://github.com/Commutator-IO/learn-lean/blob/main/courses/01-college/01-nombres-et-calculs/EcrituresDesNombres.lean\#L13}{EcrituresDesNombres.lean\#L13}

\begin{theoreme}
Le chiffre de rang \texttt{i} s'obtient en supprimant les \texttt{i} premiers chiffres, puis en
prenant le reste modulo dix.
\end{theoreme}

\begin{proof}
Diviser par $10^i$ efface les $i$ derniers chiffres, et prendre le reste modulo dix
ne garde que le dernier de ceux qui restent. Le membre de gauche fait la même chose dans
l'ordre inverse : on tronque d'abord aux $i+1$ derniers chiffres, puis on divise par $10^i$.
Les deux lectures coïncident parce que $10^{i+1} = 10^i \times 10$.
\end{proof}
\source{https://github.com/Commutator-IO/learn-lean/blob/main/courses/01-college/01-nombres-et-calculs/EcrituresDesNombres.lean\#L19}{EcrituresDesNombres.lean\#L19}

\subsection{Ordre sur les décimaux ; comparaison, encadrement, intercalation}

\begin{theoreme}
Entre deux décimaux distincts s'intercale toujours un troisième nombre.
\end{theoreme}

\begin{proof}
Les rationnels sont denses : entre deux d'entre eux s'en trouve toujours un
troisième, par exemple leur moyenne.
\end{proof}
\source{https://github.com/Commutator-IO/learn-lean/blob/main/courses/01-college/01-nombres-et-calculs/EcrituresDesNombres.lean\#L26}{EcrituresDesNombres.lean\#L26}

\begin{theoreme}
Encadrement d'un rationnel par deux entiers consécutifs.
\end{theoreme}

\begin{proof}
La partie entière $\lfloor x \rfloor$ est le plus grand entier inférieur ou égal à
$x$ : elle vérifie donc $\lfloor x \rfloor \le x$, et $x < \lfloor x \rfloor + 1$ puisque
l'entier suivant est strictement plus grand que $x$.
\end{proof}
\source{https://github.com/Commutator-IO/learn-lean/blob/main/courses/01-college/01-nombres-et-calculs/EcrituresDesNombres.lean\#L30}{EcrituresDesNombres.lean\#L30}

\subsection{Ordre sur les relatifs ; opposé, distance à zéro}

\begin{theoreme}
Prendre l'opposé renverse l'ordre.
\end{theoreme}

\begin{proof}
Prendre l'opposé est une bijection décroissante : $a < b$ équivaut à $-b < -a$.
\end{proof}
\source{https://github.com/Commutator-IO/learn-lean/blob/main/courses/01-college/01-nombres-et-calculs/EcrituresDesNombres.lean\#L36}{EcrituresDesNombres.lean\#L36}

\begin{theoreme}
La distance à zéro est le plus grand des deux nombres opposés.
\end{theoreme}

\begin{proof}
La distance à zéro vaut $a$ si $a$ est positif, $-a$ sinon : dans les deux cas,
c'est le plus grand des deux nombres $a$ et $-a$.
\end{proof}
\source{https://github.com/Commutator-IO/learn-lean/blob/main/courses/01-college/01-nombres-et-calculs/EcrituresDesNombres.lean\#L39}{EcrituresDesNombres.lean\#L39}

\subsection{Égalité de fractions : \texttt{a/b = (ka)/(kb)} pour \texttt{k \(\ne\) 0}}

\begin{theoreme}
Multiplier numérateur et dénominateur par un même nombre non nul ne change pas la
fraction.
\end{theoreme}

\begin{proof}
Simplifier le facteur commun $k$, non nul, au numérateur et au dénominateur.
\end{proof}
\source{https://github.com/Commutator-IO/learn-lean/blob/main/courses/01-college/01-nombres-et-calculs/EcrituresDesNombres.lean\#L45}{EcrituresDesNombres.lean\#L45}

\subsection{Comparaison de fractions ; mise au même dénominateur}

\begin{theoreme}
Deux fractions de dénominateurs positifs se comparent par produits croisés.
\end{theoreme}

\begin{proof}
Les deux dénominateurs étant strictement positifs, multiplier l'inégalité par $bd$
la conserve : $a/b < c/d$ équivaut à $ad < cb$.
\end{proof}
\source{https://github.com/Commutator-IO/learn-lean/blob/main/courses/01-college/01-nombres-et-calculs/EcrituresDesNombres.lean\#L51}{EcrituresDesNombres.lean\#L51}

\begin{theoreme}
Mise au même dénominateur.
\end{theoreme}

\begin{proof}
On réduit les deux fractions au dénominateur commun $bd$, non nul.
\end{proof}
\source{https://github.com/Commutator-IO/learn-lean/blob/main/courses/01-college/01-nombres-et-calculs/EcrituresDesNombres.lean\#L55}{EcrituresDesNombres.lean\#L55}

\subsection{Une fraction n'a pas toujours d'écriture décimale exacte (\texttt{1/3})}

\begin{theoreme}
\texttt{1/3} ne s'écrit avec aucun nombre fini de décimales : sinon \texttt{3} diviserait une
puissance de dix.
\end{theoreme}

\begin{proof}
Par l'absurde : si $1/3$ s'écrivait avec un nombre fini de décimales, on aurait
$10^n = 3a$ pour un entier $a$, donc $3$ diviserait $10^n$. Or $3$ est premier, il devrait
donc diviser $10$, ce qui est faux.
\end{proof}
\source{https://github.com/Commutator-IO/learn-lean/blob/main/courses/01-college/01-nombres-et-calculs/EcrituresDesNombres.lean\#L63}{EcrituresDesNombres.lean\#L63}

\subsection{Arrondi, troncature, valeur approchée à \texttt{10⁻ⁿ} près, encadrement}

\begin{theoreme}
La troncature à l'unité est le plancher : elle encadre le nombre à moins de un.
\end{theoreme}

\begin{proof}
La troncature à l'unité est la partie entière : elle est inférieure ou égale à $x$,
et $x$ lui est inférieur de moins de un puisque $x < \lfloor x \rfloor + 1$.
\end{proof}
\source{https://github.com/Commutator-IO/learn-lean/blob/main/courses/01-college/01-nombres-et-calculs/EcrituresDesNombres.lean\#L72}{EcrituresDesNombres.lean\#L72}

\begin{theoreme}
L'arrondi à \texttt{10⁻ⁿ} près approche le nombre à \texttt{10⁻ⁿ / 2} près.
\end{theoreme}

\begin{proof}
L'arrondi d'un nombre en est distant d'au plus $1/2$. En l'appliquant à
$x \times 10^n$ puis en divisant par $10^n$, l'écart est divisé par $10^n$ : il vaut au plus
$1/(2 \times 10^n)$.
\end{proof}
\source{https://github.com/Commutator-IO/learn-lean/blob/main/courses/01-college/01-nombres-et-calculs/EcrituresDesNombres.lean\#L76}{EcrituresDesNombres.lean\#L76}

\subsection{Ordre de grandeur d'un résultat ; contrôle de la vraisemblance d'un calcul}

\begin{theoreme}
Le produit de deux nombres d'ordres de grandeur donnés a son ordre de grandeur
encadré : c'est ce qui permet de contrôler un calcul à vue.
\end{theoreme}

\begin{proof}
Le produit des minorants minore le produit, et le produit des majorants le majore,
les quatre nombres étant positifs :
\[
10^{p+q} = 10^p \times 10^q \le xy < 10^{p+1} \times 10^{q+1} = 10^{p+q+2} .
\]
L'ordre de grandeur du produit est donc connu à un facteur cent près, ce qui suffit à
contrôler un calcul à vue.
\end{proof}
\source{https://github.com/Commutator-IO/learn-lean/blob/main/courses/01-college/01-nombres-et-calculs/EcrituresDesNombres.lean\#L90}{EcrituresDesNombres.lean\#L90}

\subsection{Écriture scientifique : existence et unicité de \texttt{a \(\times\) 10ⁿ} avec \texttt{1 \(\le\) |a| < 10}}

\begin{theoreme}
Tout réel non nul s'écrit \texttt{a \(\times\) 10ⁿ} avec \texttt{1 \(\le\) |a| < 10}.
\end{theoreme}

\begin{proof}
On prend pour exposant $n$ le logarithme entier de $|x|$ en base dix, c'est-à-dire
le plus grand entier tel que $10^n \le |x|$, et pour mantisse $a = x / 10^n$. L'encadrement
$10^n \le |x| < 10^{n+1}$ donne alors $1 \le |a| < 10$.
\end{proof}
\source{https://github.com/Commutator-IO/learn-lean/blob/main/courses/01-college/01-nombres-et-calculs/EcrituresDesNombres.lean\#L105}{EcrituresDesNombres.lean\#L105}

\begin{theoreme}
Unicité de l'exposant, donc de l'écriture.
\end{theoreme}

\begin{proof}
Le rapport des deux écritures est une puissance de dix : $10^{n-n'} = |a'|/|a|$.
Les deux mantisses étant dans $[1, 10[$, ce rapport est strictement compris entre $10^{-1}$
et $10^{1}$, donc son exposant vérifie $-1 < n - n' < 1$ : il est nul. Les exposants étant
égaux, la simplification par $10^n$ donne $a = a'$.
\end{proof}
\source{https://github.com/Commutator-IO/learn-lean/blob/main/courses/01-college/01-nombres-et-calculs/EcrituresDesNombres.lean\#L120}{EcrituresDesNombres.lean\#L120}

\subsection{Racine carrée : \texttt{(\(\sqrt{a}\))² = a} et \texttt{\(\sqrt{a^{2}}\) = a} pour \texttt{a \(\ge\) 0}}

\begin{theoreme}
Élever au carré annule la racine.
\end{theoreme}

\begin{proof}
C'est la définition de la racine carrée d'un nombre positif : c'est l'unique
nombre positif dont le carré vaut $a$.
\end{proof}
\source{https://github.com/Commutator-IO/learn-lean/blob/main/courses/01-college/01-nombres-et-calculs/EcrituresDesNombres.lean\#L163}{EcrituresDesNombres.lean\#L163}

\begin{theoreme}
Et réciproquement, pour un nombre positif.
\end{theoreme}

\begin{proof}
Réciproquement, la racine du carré d'un nombre positif redonne ce nombre. Pour un
nombre négatif, elle en donnerait la valeur absolue.
\end{proof}
\source{https://github.com/Commutator-IO/learn-lean/blob/main/courses/01-college/01-nombres-et-calculs/EcrituresDesNombres.lean\#L167}{EcrituresDesNombres.lean\#L167}

\subsection{\texttt{\(\sqrt{ab}\) = \(\sqrt{a}\) \(\times\) \(\sqrt{b}\)} et \texttt{\(\sqrt{a/b}\) = \(\sqrt{a}\) / \(\sqrt{b}\)} (\texttt{a \(\ge\) 0}, \texttt{b > 0})}

\begin{theoreme}
La racine d'un produit est le produit des racines.
\end{theoreme}

\begin{proof}
La racine du produit et le produit des racines sont deux nombres positifs de même
carré, à savoir $ab$ : ils sont donc égaux.
\end{proof}
\source{https://github.com/Commutator-IO/learn-lean/blob/main/courses/01-college/01-nombres-et-calculs/EcrituresDesNombres.lean\#L173}{EcrituresDesNombres.lean\#L173}

\begin{theoreme}
La racine d'un quotient est le quotient des racines.
\end{theoreme}

\begin{proof}
Même argument avec un quotient.
\end{proof}
\source{https://github.com/Commutator-IO/learn-lean/blob/main/courses/01-college/01-nombres-et-calculs/EcrituresDesNombres.lean\#L177}{EcrituresDesNombres.lean\#L177}

\subsection{Contre-exemple : \texttt{\(\sqrt{a + b}\) \(\ne\) \(\sqrt{a}\) + \(\sqrt{b}\)} en général}

\begin{exemple}
La racine d'une somme n'est pas la somme des racines : \texttt{\(\sqrt{2}\) \(\ne\) 2}.
\end{exemple}

\begin{explication}
Si l'on avait $\sqrt{1+1} = \sqrt1 + \sqrt1 = 2$, alors en élevant au carré on
obtiendrait $2 = 4$. La racine d'une somme n'est donc pas la somme des racines.
\end{explication}
\source{https://github.com/Commutator-IO/learn-lean/blob/main/courses/01-college/01-nombres-et-calculs/EcrituresDesNombres.lean\#L183}{EcrituresDesNombres.lean\#L183}

\subsection{Rationnels et irrationnels : \texttt{\(\sqrt{2}\)} n'est pas rationnel}

\begin{theoreme}
\texttt{\(\sqrt{2}\)} est irrationnel. Au collège, le résultat est mentionné sans démonstration.
\end{theoreme}

\begin{proof}
Mathlib ne procède pas par la parité, comme le fait le raisonnement scolaire, mais par une
forme du théorème des racines rationnelles.

Supposons $\sqrt 2 = N/D$, la fraction étant écrite sous forme réduite. En élevant au carré,
$N^2 = 2D^2$, donc $D^2$ divise $N^2$. Or, pour des entiers, $D^n \mid N^n$ entraîne
$D \mid N$ ; et comme $N$ et $D$ sont premiers entre eux, il vient $D = 1$. Ainsi $\sqrt 2$
serait un entier — or aucun entier n'a $2$ pour carré, puisque $1^2 = 1$ et $2^2 = 4$.

L'argument ne dépend pas de $2$ : il vaut pour $\sqrt p$ avec $p$ premier, et plus
généralement pour toute racine $n$-ième d'un entier qui n'est pas une puissance $n$-ième
exacte (\texttt{irrational\_nrt\_of\_notint\_nrt}).

Au collège, l'irrationalité de $\sqrt 2$ est mentionnée sans démonstration.
\end{proof}
\source{https://github.com/Commutator-IO/learn-lean/blob/main/courses/01-college/01-nombres-et-calculs/EcrituresDesNombres.lean\#L194}{EcrituresDesNombres.lean\#L194}

\section{EntiersDivisibilite}

\noindent
Collège \textemdash{} section \og{} Entiers, divisibilité \fg{}.
Les preuves n'utilisent que la bibliothèque standard ; Mathlib n'est importé que pour
ses notations (\(\mathbb{N}\)) et la cohérence avec les autres fichiers du chapitre.
Énoncés et démonstrations en français : voir NombresEtCalculs.tex.

\subsection{Définitions}

\noindent
Toutes les définitions du fichier sont posées ici, avant les démonstrations.

\begin{definition}
\texttt{n} est \emph{pair} lorsqu'il est le double d'un entier naturel :
\[ n \text{ pair} \iff \exists\, k \in \mathbb{N},\ n = 2k . \]
\end{definition}
\source{https://github.com/Commutator-IO/learn-lean/blob/main/courses/01-college/01-nombres-et-calculs/EntiersDivisibilite.lean\#L16}{EntiersDivisibilite.lean\#L16}

\begin{definition}
\texttt{n} est \emph{impair} lorsqu'il est le double d'un entier naturel, plus un :
\[ n \text{ impair} \iff \exists\, k \in \mathbb{N},\ n = 2k + 1 . \]
\end{definition}
\source{https://github.com/Commutator-IO/learn-lean/blob/main/courses/01-college/01-nombres-et-calculs/EntiersDivisibilite.lean\#L19}{EntiersDivisibilite.lean\#L19}

\begin{definition}
\texttt{p} est \emph{premier} lorsqu'il vaut au moins deux et n'admet pas d'autre
diviseur que un et lui-même :
\[ p \text{ premier} \iff 2 \le p \ \text{ et } \ \forall d,\ d \mid p \Rightarrow
   (d = 1 \ \text{ ou } \ d = p) . \]
\end{definition}
\source{https://github.com/Commutator-IO/learn-lean/blob/main/courses/01-college/01-nombres-et-calculs/EntiersDivisibilite.lean\#L22}{EntiersDivisibilite.lean\#L22}

\begin{definition}
Le \emph{chiffre des unités} de $n$ est le reste de sa division par dix :
\[ u(n) = n \bmod 10 . \]
\end{definition}
\source{https://github.com/Commutator-IO/learn-lean/blob/main/courses/01-college/01-nombres-et-calculs/EntiersDivisibilite.lean\#L25}{EntiersDivisibilite.lean\#L25}

\begin{definition}
La \emph{somme des chiffres} se calcule par récursion, en retirant le chiffre des
unités puis en recommençant sur ce qui reste :
\[ S(0) = 0, \qquad
   S(n) = (n \bmod 10) + S\!\left(\left\lfloor \tfrac{n}{10} \right\rfloor\right)
   \ \text{ pour } n \ne 0 . \]
La récursion termine car $\lfloor n/10 \rfloor < n$ dès que $n \ne 0$.
\end{definition}
\source{https://github.com/Commutator-IO/learn-lean/blob/main/courses/01-college/01-nombres-et-calculs/EntiersDivisibilite.lean\#L28}{EntiersDivisibilite.lean\#L28}

\begin{definition}
Le nombre formé par les \emph{deux derniers chiffres} de $n$ est le reste de sa
division par cent :
\[ d(n) = n \bmod 100 . \]
\end{definition}
\source{https://github.com/Commutator-IO/learn-lean/blob/main/courses/01-college/01-nombres-et-calculs/EntiersDivisibilite.lean\#L34}{EntiersDivisibilite.lean\#L34}

\begin{definition}
Le \emph{crible} s'écrit comme un test calculable : $n$ passe le test lorsqu'il vaut au
moins deux et qu'aucun entier compris entre deux et $n-1$ ne le divise :
\[ \textsf{estPremier}(n) = \bigl(2 \le n\bigr) \ \text{ et } \
   \forall d < n,\ \bigl(d < 2 \ \text{ ou } \ n \bmod d \ne 0\bigr) . \]
\end{definition}
\source{https://github.com/Commutator-IO/learn-lean/blob/main/courses/01-college/01-nombres-et-calculs/EntiersDivisibilite.lean\#L37}{EntiersDivisibilite.lean\#L37}

\begin{definition}
Le \emph{produit} d'une liste se calcule par récursion sur celle-ci :
\[ \Pi(\,[\,]\,) = 1, \qquad \Pi(p :: \ell) = p \times \Pi(\ell) . \]
\end{definition}
\source{https://github.com/Commutator-IO/learn-lean/blob/main/courses/01-college/01-nombres-et-calculs/EntiersDivisibilite.lean\#L41}{EntiersDivisibilite.lean\#L41}

\subsection{Un entier est pair ou impair, jamais les deux}

\begin{theoreme}
Un entier est pair ou impair.
\end{theoreme}

\begin{proof}
Le reste $n \bmod 2$ ne peut valoir que $0$ ou $1$ ; distinguons ces deux cas. Dans
l'un comme dans l'autre, il suffit de prendre $k = \lfloor n/2 \rfloor$ : de
$n = 2\lfloor n/2 \rfloor + (n \bmod 2)$ on tire $n = 2k$ si le reste est nul, et
$n = 2k + 1$ sinon. Les deux égalités se vérifient par un calcul.
\end{proof}
\source{https://github.com/Commutator-IO/learn-lean/blob/main/courses/01-college/01-nombres-et-calculs/EntiersDivisibilite.lean\#L48}{EntiersDivisibilite.lean\#L48}

\begin{theoreme}
Aucun entier n'est à la fois pair et impair.
\end{theoreme}

\begin{proof}
Par l'absurde : si $n$ était à la fois pair et impair, on écrirait $n = 2k$ et
$n = 2l + 1$, d'où $2k = 2l + 1$. Or un nombre pair ne peut pas être égal à un nombre
impair.
\end{proof}
\source{https://github.com/Commutator-IO/learn-lean/blob/main/courses/01-college/01-nombres-et-calculs/EntiersDivisibilite.lean\#L54}{EntiersDivisibilite.lean\#L54}

\begin{theoreme}
Réunion des deux énoncés précédents : tout entier relève d'un cas, et d'un seul.
\end{theoreme}

\begin{proof}
C'est la conjonction des deux théorèmes précédents.
\end{proof}
\source{https://github.com/Commutator-IO/learn-lean/blob/main/courses/01-college/01-nombres-et-calculs/EntiersDivisibilite.lean\#L59}{EntiersDivisibilite.lean\#L59}

\subsection{Somme de deux pairs = pair ; pair + impair = impair ; impair + impair = pair}

\begin{theoreme}
Somme de deux pairs = pair.
\end{theoreme}

\begin{proof}
On écrit les hypothèses sous la forme $m = 2k$ et $n = 2l$. Alors
$m + n = 2k + 2l = 2(k+l)$, et $k + l$ convient.
\end{proof}
\source{https://github.com/Commutator-IO/learn-lean/blob/main/courses/01-college/01-nombres-et-calculs/EntiersDivisibilite.lean\#L66}{EntiersDivisibilite.lean\#L66}

\begin{theoreme}
Pair + impair = impair.
\end{theoreme}

\begin{proof}
On écrit $m = 2k$ et $n = 2l + 1$. Alors $m + n = 2(k+l) + 1$ : on prend encore
$k + l$, et le $+1$ de $n$ devient celui de la conclusion.
\end{proof}
\source{https://github.com/Commutator-IO/learn-lean/blob/main/courses/01-college/01-nombres-et-calculs/EntiersDivisibilite.lean\#L73}{EntiersDivisibilite.lean\#L73}

\begin{theoreme}
Impair + impair = pair.
\end{theoreme}

\begin{proof}
On écrit $m = 2k + 1$ et $n = 2l + 1$. Alors $m + n = 2k + 2l + 2 = 2(k + l + 1)$ : les
deux $+1$ se recombinent en un facteur $2$, et $k + l + 1$ convient.
\end{proof}
\source{https://github.com/Commutator-IO/learn-lean/blob/main/courses/01-college/01-nombres-et-calculs/EntiersDivisibilite.lean\#L80}{EntiersDivisibilite.lean\#L80}

\subsection{Notion de multiple et de diviseur ; un diviseur de \texttt{n} est inférieur ou égal à \texttt{n}}

\begin{theoreme}
Notion de multiple et de diviseur.
\end{theoreme}

\begin{proof}
Les deux membres disent la même chose : sur $\mathbb{N}$, « $a$ divise $n$ » signifie
précisément « il existe un entier $k$ tel que $n = ak$ ». Il n'y a rien à vérifier.
\end{proof}
\source{https://github.com/Commutator-IO/learn-lean/blob/main/courses/01-college/01-nombres-et-calculs/EntiersDivisibilite.lean\#L89}{EntiersDivisibilite.lean\#L89}

\begin{theoreme}
Un diviseur de \texttt{n} est inférieur ou égal à \texttt{n}.
\end{theoreme}

\begin{proof}
Si $n = ak$ avec $n > 0$,
alors $k \ge 1$, donc $a \le ak = n$.
\end{proof}
\source{https://github.com/Commutator-IO/learn-lean/blob/main/courses/01-college/01-nombres-et-calculs/EntiersDivisibilite.lean\#L92}{EntiersDivisibilite.lean\#L92}

\begin{theoreme}
Contre-exemple pour \texttt{n = 0}.
\end{theoreme}

\begin{proof}
Il suffit de prendre $k = 0$, puisque $0 = a \cdot 0$.

Cet énoncé est la raison d'être de l'hypothèse $0 < n$ du théorème précédent : les
diviseurs de $0$ ne sont pas bornés. L'énoncé de la classe, « un diviseur de $n$ est
inférieur ou égal à $n$ », est donc faux tel quel ; il n'est vrai que parce qu'au
collège $n$ compte toujours des objets à partager.
\end{proof}
\source{https://github.com/Commutator-IO/learn-lean/blob/main/courses/01-college/01-nombres-et-calculs/EntiersDivisibilite.lean\#L96}{EntiersDivisibilite.lean\#L96}

\subsection{Critère de divisibilité par 2, par 5, par 10 (chiffre des unités)}

\begin{theoreme}
Critère de divisibilité par 2 (chiffre des unités).
\end{theoreme}

\begin{proof}
Après dépliage de la définition du chiffre des unités, il s'agit de
$2 \mid n \iff 2 \mid (n \bmod 10)$, prouvée dans les deux sens : comme $2 \mid 10$, écrire $n = 10q + (n \bmod 10)$ montre que $n$ et
$n \bmod 10$ ont le même reste modulo $2$.
\end{proof}
\source{https://github.com/Commutator-IO/learn-lean/blob/main/courses/01-college/01-nombres-et-calculs/EntiersDivisibilite.lean\#L102}{EntiersDivisibilite.lean\#L102}

\begin{theoreme}
Critère de divisibilité par 2, chiffres énumérés.
\end{theoreme}

\begin{proof}
Même argument, la condition $2 \mid (n \bmod 10)$ étant énumérée : le reste modulo $10$
est l'un des dix chiffres, et les pairs sont $0$, $2$, $4$, $6$, $8$.
\end{proof}
\source{https://github.com/Commutator-IO/learn-lean/blob/main/courses/01-college/01-nombres-et-calculs/EntiersDivisibilite.lean\#L107}{EntiersDivisibilite.lean\#L107}

\begin{theoreme}
Critère de divisibilité par 5 (chiffre des unités).
\end{theoreme}

\begin{proof}
Identique, avec $5 \mid 10$ : $n$ et $n \bmod 10$ ont même reste modulo $5$, et les
chiffres divisibles par $5$ sont $0$ et $5$.
\end{proof}
\source{https://github.com/Commutator-IO/learn-lean/blob/main/courses/01-college/01-nombres-et-calculs/EntiersDivisibilite.lean\#L114}{EntiersDivisibilite.lean\#L114}

\begin{theoreme}
Critère de divisibilité par 10 (chiffre des unités).
\end{theoreme}

\begin{proof}
Cas limite du même argument : $n \bmod 10$ est le reste de $n$ modulo $10$, donc il est
nul exactement quand $10$ divise $n$.
\end{proof}
\source{https://github.com/Commutator-IO/learn-lean/blob/main/courses/01-college/01-nombres-et-calculs/EntiersDivisibilite.lean\#L120}{EntiersDivisibilite.lean\#L120}

\subsection{Critère de divisibilité par 3 et par 9 (somme des chiffres)}

\begin{theoreme}
\texttt{n} et la somme de ses chiffres ont le même reste modulo 9.
\end{theoreme}

\begin{proof}
Par récurrence forte sur $n$. Si $n = 0$, les deux membres valent $0$.

Sinon, on déplie $S(n) = (n \bmod 10) + S(\lfloor n/10 \rfloor)$. Le quotient
$\lfloor n/10 \rfloor$ est strictement plus petit que $n$, ce qui autorise l'hypothèse de
récurrence : $\lfloor n/10 \rfloor \equiv S(\lfloor n/10 \rfloor) \pmod 9$. Comme
$n = 10\lfloor n/10 \rfloor + (n \bmod 10)$ et $10 \equiv 1 \pmod 9$, on obtient
\[
n \equiv \lfloor n/10 \rfloor + (n \bmod 10)
  \equiv S(\lfloor n/10 \rfloor) + (n \bmod 10) = S(n) \pmod 9 .
\]
\end{proof}
\source{https://github.com/Commutator-IO/learn-lean/blob/main/courses/01-college/01-nombres-et-calculs/EntiersDivisibilite.lean\#L127}{EntiersDivisibilite.lean\#L127}

\begin{theoreme}
\texttt{n} et la somme de ses chiffres ont le même reste modulo 3.
\end{theoreme}

\begin{proof}
Même récurrence, en remplaçant $9$ par $3$ : $10 \equiv 1 \pmod 3$ également.
\end{proof}
\source{https://github.com/Commutator-IO/learn-lean/blob/main/courses/01-college/01-nombres-et-calculs/EntiersDivisibilite.lean\#L139}{EntiersDivisibilite.lean\#L139}

\begin{theoreme}
Critère de divisibilité par 3 (somme des chiffres).
\end{theoreme}

\begin{proof}
Conséquence immédiate du lemme : $n$ et $S(n)$ ont le même reste modulo $3$, donc l'un
est divisible par $3$ si et seulement si l'autre l'est.
\end{proof}
\source{https://github.com/Commutator-IO/learn-lean/blob/main/courses/01-college/01-nombres-et-calculs/EntiersDivisibilite.lean\#L151}{EntiersDivisibilite.lean\#L151}

\begin{theoreme}
Critère de divisibilité par 9 (somme des chiffres).
\end{theoreme}

\begin{proof}
Identique avec le lemme modulo $9$.
\end{proof}
\source{https://github.com/Commutator-IO/learn-lean/blob/main/courses/01-college/01-nombres-et-calculs/EntiersDivisibilite.lean\#L156}{EntiersDivisibilite.lean\#L156}

\subsection{Critère de divisibilité par 4 (deux derniers chiffres)}

\begin{theoreme}
Critère de divisibilité par 4 (deux derniers chiffres).
\end{theoreme}

\begin{proof}
Même schéma que pour $2$ et $5$, cette fois avec $4 \mid 100$ : écrire
$n = 100q + (n \bmod 100)$ montre que $n$ et $n \bmod 100$ ont même reste modulo $4$.
\end{proof}
\source{https://github.com/Commutator-IO/learn-lean/blob/main/courses/01-college/01-nombres-et-calculs/EntiersDivisibilite.lean\#L163}{EntiersDivisibilite.lean\#L163}

\subsection{Si \texttt{a \(\mid\) b} et \texttt{a \(\mid\) c} alors \texttt{a \(\mid\) (b + c)} et \texttt{a \(\mid\) (b − c)}}

\begin{theoreme}
Si \texttt{a \(\mid\) b} et \texttt{a \(\mid\) c} alors \texttt{a \(\mid\) (b + c)}.
\end{theoreme}

\begin{proof}
Si $b = ak$ et $c = al$, alors $b + c = a(k+l)$ (la stabilité de la divisibilité par somme).
\end{proof}
\source{https://github.com/Commutator-IO/learn-lean/blob/main/courses/01-college/01-nombres-et-calculs/EntiersDivisibilite.lean\#L170}{EntiersDivisibilite.lean\#L170}

\begin{theoreme}
Si \texttt{a \(\mid\) b} et \texttt{a \(\mid\) c} alors \texttt{a \(\mid\) (b − c)}.
\end{theoreme}

\begin{proof}
Si $b = ak$ et $c = al$, alors $b - c = a(k - l)$ : un diviseur commun divise la
différence.

L'énoncé mathématique suppose $c \le b$, faute de quoi $b - c$ n'est pas un entier naturel.
Le théorème Lean s'en passe, mais par convention : sur $\mathbb{N}$ la soustraction est
tronquée, donc $b - c$ vaut $0$ dès que $b < c$, et tout entier divise $0$. L'égalité est
alors vraie sans rien dire du cas visé par l'énoncé.
\end{proof}
\source{https://github.com/Commutator-IO/learn-lean/blob/main/courses/01-college/01-nombres-et-calculs/EntiersDivisibilite.lean\#L174}{EntiersDivisibilite.lean\#L174}

\subsection{Division euclidienne : existence et unicité de \texttt{(q, r)} avec \texttt{a = bq + r}, \texttt{0 \(\le\) r < b}}

\begin{theoreme}
Division euclidienne : existence de \texttt{(q, r)} avec \texttt{a = bq + r}, \texttt{0 \(\le\) r < b}.
\end{theoreme}

\begin{proof}
Prendre $q = \lfloor a/b \rfloor$ et $r = a \bmod b$. L'égalité $a = bq + r$ est
l'égalité de la division euclidienne, et l'inégalité $r < b$ est la majoration du reste, qui demande
$b > 0$.
\end{proof}
\source{https://github.com/Commutator-IO/learn-lean/blob/main/courses/01-college/01-nombres-et-calculs/EntiersDivisibilite.lean\#L180}{EntiersDivisibilite.lean\#L180}

\begin{theoreme}
Division euclidienne : unicité de \texttt{(q, r)}.
\end{theoreme}

\begin{proof}
Le critère le critère d'unicité du quotient et du reste dit que, pour $b > 0$, la conjonction de
$r + bq = a$ et $r < b$ caractérise $q = \lfloor a/b \rfloor$ et $r = a \bmod b$.

La première écriture $a = bq + r$ avec $r < b$ donne donc $q = \lfloor a/b \rfloor$ et
$r = a \bmod b$ ; la seconde donne de même $q' = \lfloor a/b \rfloor$ et
$r' = a \bmod b$. D'où $q = q'$ et $r = r'$.
\end{proof}
\source{https://github.com/Commutator-IO/learn-lean/blob/main/courses/01-college/01-nombres-et-calculs/EntiersDivisibilite.lean\#L185}{EntiersDivisibilite.lean\#L185}

\subsection{Nombre premier ; tout entier > 1 admet un diviseur premier}

\begin{theoreme}
Nombre premier : ses seuls diviseurs sont 1 et lui-même.
\end{theoreme}

\begin{proof}
C'est la seconde composante de la définition de \texttt{Premier}.
\end{proof}
\source{https://github.com/Commutator-IO/learn-lean/blob/main/courses/01-college/01-nombres-et-calculs/EntiersDivisibilite.lean\#L197}{EntiersDivisibilite.lean\#L197}

\begin{theoreme}
Tout entier > 1 admet un diviseur premier.
\end{theoreme}

\begin{proof}
Par récurrence forte sur $n$. Distinguons deux cas.

Si $n$ est premier, il est son propre diviseur premier.

Sinon, la négation de la définition fournit, par l'absurde, un diviseur $d$ de $n$ avec
$d \ne 1$ et $d \ne n$. Ce $d$ n'est pas nul : $0 \mid n$ entraînerait $n = 0$, contraire
à $n \ge 2$. Donc $2 \le d$. De $d \mid n$ et $n > 0$ on tire $d \le n$, et $d \ne n$
donne $d < n$. L'hypothèse de récurrence appliquée à $d$ fournit un premier $p$ divisant
$d$, et $p \mid n$ par transitivité.
\end{proof}
\source{https://github.com/Commutator-IO/learn-lean/blob/main/courses/01-college/01-nombres-et-calculs/EntiersDivisibilite.lean\#L201}{EntiersDivisibilite.lean\#L201}

\subsection{Crible d'Ératosthène : lister les nombres premiers inférieurs à 100}

\begin{theoreme}
Le test et la définition coïncident.
\end{theoreme}

\begin{proof}
($\Rightarrow$) Supposons le test satisfait et soit $d \mid n$. Comme $n \ge 2 > 0$, on a
$d \le n$. Si $d = n$, la conclusion est acquise. Si $d < n$, le test donne $d < 2$ ou
$n \bmod d \ne 0$ ; or $d \mid n$ force $n \bmod d = 0$, donc $d < 2$, c'est-à-dire
$d = 0$ ou $d = 1$. Le cas $d = 0$ est exclu, car $0 \mid n$ donnerait $n = 0$. Reste
$d = 1$.

($\Leftarrow$) Supposons $n$ premier et soit $d < n$. Si $d < 2$, la disjonction du test
est satisfaite. Sinon, supposons $n \bmod d = 0$ : alors $d \mid n$, donc $d = 1$ ou
$d = n$ par primalité, ce qui contredit $2 \le d < n$.
\end{proof}
\source{https://github.com/Commutator-IO/learn-lean/blob/main/courses/01-college/01-nombres-et-calculs/EntiersDivisibilite.lean\#L227}{EntiersDivisibilite.lean\#L227}

\begin{exemple}
Crible d'Ératosthène : les nombres premiers inférieurs à 100.
\end{exemple}

\begin{explication}
Par calcul : le noyau évalue la liste $0, 1, \dots, 99$, la filtre par le test et compare
le résultat à la liste annoncée. Le crible est donc exécuté, et pas seulement décrit —
c'est le théorème précédent qui garantit que ce que le test retient est bien la
primalité.
\end{explication}
\source{https://github.com/Commutator-IO/learn-lean/blob/main/courses/01-college/01-nombres-et-calculs/EntiersDivisibilite.lean\#L250}{EntiersDivisibilite.lean\#L250}

\subsection{Décomposition en produit de facteurs premiers}

\begin{theoreme}
Décomposition en produit de facteurs premiers.
\end{theoreme}

\begin{proof}
Par récurrence forte sur $n$. Si $n = 1$, la liste vide convient, son produit valant $1$.

Sinon $n \ge 2$, et le théorème précédent fournit un premier $p$ divisant $n$, disons
$n = pm$. Alors $m \ne 0$, sinon $n$ serait nul ; et $m < n$, car
$m = 1 \cdot m < p \cdot m = n$ puisque $p \ge 2$ et $m > 0$. L'hypothèse de récurrence
donne une liste $l$ de premiers de produit $m$ ; la liste $p :: l$ convient alors pour
$n$ : tous ses éléments sont premiers, et son produit vaut $p \cdot m = n$.
\end{proof}
\source{https://github.com/Commutator-IO/learn-lean/blob/main/courses/01-college/01-nombres-et-calculs/EntiersDivisibilite.lean\#L259}{EntiersDivisibilite.lean\#L259}

\begin{lemme}
Lemme d'Euclide, sur lequel repose l'unicité.
\end{lemme}

\begin{proof}
Distinguons selon que $p$ divise $a$ ou non. S'il le divise, c'est fini.

Sinon, le pgcd $\gcd(p, a)$ divise $p$, donc vaut $1$ ou $p$ par primalité de $p$. S'il
valait $p$, alors $p$ diviserait $a$, ce qui est exclu. Donc $\gcd(p,a) = 1$ : $p$ et $a$
sont premiers entre eux, et de $p \mid ab$ on conclut $p \mid b$
(le lemme de Gauss).
\end{proof}
\source{https://github.com/Commutator-IO/learn-lean/blob/main/courses/01-college/01-nombres-et-calculs/EntiersDivisibilite.lean\#L285}{EntiersDivisibilite.lean\#L285}

\begin{theoreme}
Le prédicat \texttt{Premier} défini ici coïncide avec celui de la bibliothèque.
\end{theoreme}

\begin{proof}
Les deux définitions disent la même chose, mot pour mot : un entier est premier lorsqu'il
vaut au moins $2$ et que tout diviseur de cet entier est $1$ ou lui-même. Il n'y a donc
rien à démontrer, seulement à constater que les deux énoncés se superposent.
\end{proof}
\source{https://github.com/Commutator-IO/learn-lean/blob/main/courses/01-college/01-nombres-et-calculs/EntiersDivisibilite.lean\#L297}{EntiersDivisibilite.lean\#L297}

\begin{theoreme}
Le produit défini par récursion est celui de la bibliothèque.
\end{theoreme}

\begin{proof}
Par récurrence sur la liste : le produit de la liste vide vaut un dans les deux
définitions, et l'étape de récurrence multiplie de part et d'autre par la tête de liste.
\end{proof}
\source{https://github.com/Commutator-IO/learn-lean/blob/main/courses/01-college/01-nombres-et-calculs/EntiersDivisibilite.lean\#L300}{EntiersDivisibilite.lean\#L300}

\begin{theoreme}
Unicité de la décomposition, à l'ordre des facteurs près : deux listes de nombres
premiers de même produit sont permutées l'une de l'autre.
\end{theoreme}

\begin{proof}
Les deux listes ne sont pas comparées l'une à l'autre, mais chacune à une même troisième :
celle des facteurs premiers de leur produit commun. Il suffit donc de savoir que toute liste
de nombres premiers de produit $n$ est une permutation de celle-là — c'est le théorème
fondamental de l'arithmétique, et sa démonstration est une récurrence sur la longueur de la
liste.

Soient $a :: l_1$ et $l_2$ deux listes de nombres premiers de même produit. Le nombre $a$
est premier et divise le produit de $l_2$ ; par le lemme d'Euclide, démontré plus haut, il
divise l'un de ses facteurs, lequel est premier : $a$ est donc l'un des éléments de $l_2$.
On le retire des deux listes. Les produits restants sont encore égaux, puisqu'on a simplifié
par $a$, qui n'est pas nul. L'hypothèse de récurrence s'applique aux deux restes, et
remettre $a$ en tête conclut.

Les deux cas de base disent qu'une liste vide ne saurait avoir le même produit qu'une liste
non vide de nombres premiers : $1$ n'est divisible par aucun nombre premier.

Au collège, cette unicité est admise sans même être énoncée.
\end{proof}
\source{https://github.com/Commutator-IO/learn-lean/blob/main/courses/01-college/01-nombres-et-calculs/EntiersDivisibilite.lean\#L307}{EntiersDivisibilite.lean\#L307}

\subsection{PGCD, algorithme d'Euclide}

\begin{theoreme}
PGCD, algorithme d'Euclide : \texttt{pgcd(a, b) = pgcd(b, a mod b)}.
\end{theoreme}

\begin{proof}
C'est l'équation de récursion de le pgcd : la définition de la fonction
\emph{est} l'algorithme d'Euclide, et sa terminaison tient à la décroissance stricte de
$b \bmod a$.
\end{proof}
\source{https://github.com/Commutator-IO/learn-lean/blob/main/courses/01-college/01-nombres-et-calculs/EntiersDivisibilite.lean\#L321}{EntiersDivisibilite.lean\#L321}

\begin{theoreme}
Le PGCD est un diviseur commun.
\end{theoreme}

\begin{proof}
Conjonction de « le pgcd divise le premier » et « le pgcd divise le second ».
\end{proof}
\source{https://github.com/Commutator-IO/learn-lean/blob/main/courses/01-college/01-nombres-et-calculs/EntiersDivisibilite.lean\#L325}{EntiersDivisibilite.lean\#L325}

\begin{theoreme}
Tout diviseur commun divise le PGCD.
\end{theoreme}

\begin{proof}
« tout diviseur commun divise le pgcd ». Noter que la maximalité s'exprime ici par la divisibilité et non
par l'ordre : tout diviseur commun \emph{divise} le pgcd, ce qui est plus fort que de lui
être inférieur.
\end{proof}
\source{https://github.com/Commutator-IO/learn-lean/blob/main/courses/01-college/01-nombres-et-calculs/EntiersDivisibilite.lean\#L329}{EntiersDivisibilite.lean\#L329}

\subsection{Deux nombres sont premiers entre eux \(\iff\) \texttt{pgcd = 1}}

\begin{theoreme}
Nombres premiers entre eux \(\iff\) \texttt{pgcd = 1}.
\end{theoreme}

\begin{proof}
Vrai par définition : sur $\mathbb{N}$, « premiers entre eux » \emph{est} l'égalité du
pgcd à $1$.
\end{proof}
\source{https://github.com/Commutator-IO/learn-lean/blob/main/courses/01-college/01-nombres-et-calculs/EntiersDivisibilite.lean\#L336}{EntiersDivisibilite.lean\#L336}

\begin{theoreme}
Nombres premiers entre eux \(\iff\) seuls diviseurs communs égaux à 1.
\end{theoreme}

\begin{proof}
($\Rightarrow$) Un diviseur commun $d$ de $a$ et $b$ divise leur pgcd, qui vaut $1$ ;
donc $d = 1$.

($\Leftarrow$) Il suffit d'appliquer l'hypothèse au pgcd lui-même, qui est un diviseur
commun de $a$ et $b$.
\end{proof}
\source{https://github.com/Commutator-IO/learn-lean/blob/main/courses/01-college/01-nombres-et-calculs/EntiersDivisibilite.lean\#L340}{EntiersDivisibilite.lean\#L340}

\subsection{Toute fraction admet une écriture irréductible}

\begin{theoreme}
Toute fraction admet une écriture irréductible.
\end{theoreme}

\begin{proof}
Posons $g = \gcd(n,d)$, non nul puisque $d$ l'est, puis $n' = n/g$ et $d' = d/g$.

Alors $d' > 0$, car $g \le d$ ; et $n'$ et $d'$ sont premiers entre eux
(le fait que la division par le pgcd rend les deux nombres premiers entre eux). Enfin, de $g \cdot (n/g) = n$ et
$g \cdot (d/g) = d$ :
\[
n \cdot \frac{d}{g}
 = \Bigl(g \cdot \frac{n}{g}\Bigr) \cdot \frac{d}{g}
 = \frac{n}{g} \cdot \Bigl(g \cdot \frac{d}{g}\Bigr)
 = \frac{n}{g} \cdot d ,
\]
la deuxième égalité n'étant que la commutativité et l'associativité du produit. L'égalité
des fractions est écrite en produits croisés pour rester dans $\mathbb{N}$, où $n/d$ n'a
pas de sens.
\end{proof}
\source{https://github.com/Commutator-IO/learn-lean/blob/main/courses/01-college/01-nombres-et-calculs/EntiersDivisibilite.lean\#L353}{EntiersDivisibilite.lean\#L353}

\vfill
\noindent\rule{0.35\linewidth}{0.4pt}\par
\noindent{\footnotesize Leonardo de Moura et Sebastian Ullrich,
\emph{The Lean~4 Theorem Prover and Programming Language},
\emph{Automated Deduction — CADE~28}, Springer, 2021, p.~625--635,
\href{https://doi.org/10.1007/978-3-030-79876-5_37}{doi:10.1007/978-3-030-79876-5\_37}.\par}

\end{document}
